\chapter{Introduction}
\label{ch:introduction}

\section{Chapter Overview}
\label{sec:chapter_overview}
The domain of Human-Computer Interaction (HCI) has evolved rapidly over recent years. Researchers have consistently sought natural, intuitive, and inexpensive ways for humans to interact with computers. Traditionally, physical input devices such as mechanical keyboards, membrane keypads, and touchscreens have served as primary interfaces for text and command entry. Although physical keyboards provide tactile feedback that enables rapid touch typing, their rigid mechanical construction presents several practical limitations \cite{ReviewVirtualKeyboard2020, Lee2022Virtual}. Physical keyboards occupy permanent desk space, cannot be modified once manufactured, and are difficult to sanitize in hospital wards and public kiosks \cite{ReviewVirtualKeyboard2020}.

To address these limitations, this research presents a customizable paper-based virtual keyboard framework powered by a single monocular RGB camera, AprilTag fiducial tracking, and PyTorch deep learning. Users design arbitrary key arrangements using an interactive desktop application and print them onto standard plain paper using common desktop printers. When placed on any flat surface, the printed paper is monitored by a standard webcam. Hand motions are tracked using the Google MediaPipe Hand Landmarker, and coordinates are normalized relative to hand length. A lightweight PyTorch Long Short-Term Memory (LSTM) network detects physical surface impacts across all five fingers in real time on a commodity CPU without requiring a GPU. Upon touch confirmation, planar homography maps camera pixels to layout key boundaries to trigger keystrokes or system commands. Crucially, the physical paper layout is completely decoupled from digital actions, enabling a single printed sheet to serve multiple distinct software profiles without reprinting.

This introductory chapter establishes the foundational context of the research. Section~\ref{sec:problem_background} outlines the history and limitations of existing input hardware. Section~\ref{sec:problem_statement} formulates the general and specific research problems. Section~\ref{sec:research_question} defines the primary and secondary research questions. Section~\ref{sec:research_motivation} explains the practical motivations for the study. Section~\ref{sec:research_aim} presents the overarching research aim. Section~\ref{sec:research_text_objectives} introduces the project objectives. Section~\ref{sec:rich_picture} illustrates the operational architecture. Section~\ref{sec:resource_requirements} outlines hardware and software requirements. Section~\ref{sec:project_scope} details project scope boundaries. Section~\ref{sec:thesis_organization} outlines the dissertation structure, and Section~\ref{sec:chapter_summary} summarizes the chapter.

\section{Problem Background}
\label{sec:problem_background}
Physical keyboards have remained the standard computer input interface for over half a century. While physical keys provide immediate tactile response, their fixed mechanical design prevents users from altering key sizes, reorganizing layouts for one-handed operation, or creating custom shortcut palettes \cite{ReviewVirtualKeyboard2020}.

To overcome these constraints, researchers developed optical and camera-based virtual keyboards \cite{Zhang2001Visual, Srivastava2012RealTime, Khare2019QWERTY, Ji2018Local}. Early systems employed optical laser projectors paired with infrared receivers or Time-of-Flight (ToF) depth cameras \cite{Lee2019Virtual, Kudale2016RealTime}. While depth sensors measure distance along the optical Z-axis directly, they are expensive, consume significant power, and are rarely integrated into consumer laptops. Alternative systems attempted touch detection using standard webcams by analyzing fingertip shadows \cite{Thomas2013Camera, Posner2012Single}. However, shadow tracking fails under diffuse room lighting, multiple light sources, and dark desk surfaces \cite{Yue2014Blind}.

Recent advances in computer vision have produced efficient landmark tracking libraries, such as Google MediaPipe Hands, which extracts 21 skeletal hand joints from RGB video in real time on standard CPUs \cite{Zhang2020MediaPipe, GilMartin2023Hand, Andriyanov2025Improving}. While skeletal tracking has become robust across diverse skin tones and lighting conditions \cite{SanchezBrizuela2023Lightweight}, detecting physical surface contact using a single 2D camera remains difficult \cite{Li2021CNN, Lee2022Virtual}. Standard webcams provide no depth measurements, and variations in hand size and camera distance prevent simple position thresholds from distinguishing touching from hovering \cite{Maman2023TypeNet, Yoo2026WordLevel}.

Furthermore, existing virtual keyboard implementations hard-code static QWERTY layouts into their applications. If the paper sheet moves or the camera angle changes, fixed coordinate mappings fail completely \cite{Wang2016AprilTag2, Kallwies2020Determining, Pirchheim2011Homography}. Consequently, there is an evident need for an integrated system combining custom printable layouts, fiducial planar tracking, scale-normalized hand kinematics, and real-time deep learning touch detection on commodity CPUs.

\section{Problem Statement}
\label{sec:problem_statement}
Writing a rigorous problem statement in computing research requires converting general domain issues into a specific, structured problem. Following the IRCA framework (Ideally, Reality, Consequences, and Aim) from computing research methodology, the problem addressed in this study is structured as follows:

\subsection{Ideally}
\label{subsec:irca_ideally}
Ideally, users should be able to transform any flat, ordinary desk surface into an interactive, highly customizable, and responsive multi-finger text and command input device using low-cost commodity hardware. Such an interface should operate using standard printed paper and an everyday monocular RGB webcam, running efficiently on a standard consumer CPU without requiring dedicated GPUs, specialized 3D depth sensors, or active electronic components. Users should have the freedom to design custom layouts and bind multiple distinct action profiles to the same physical sheet.

\subsection{Reality}
\label{subsec:irca_reality}
In reality, vision-based virtual keyboards remain hindered by severe practical and algorithmic limitations:
\begin{enumerate}
    \item \textbf{Single-Finger Restriction:} Prior monocular systems track only one finger at a time (usually the index digit), because optical shadow methods collapse when multiple fingers move closely together \cite{Thomas2013Camera, Posner2012Single, Ji2018Local}.
    \item \textbf{High Latency and GPU Dependence:} Modern computer vision and deep learning pipelines require heavy GPU acceleration or drop below usable frame rates on standard CPUs \cite{Enkhbat2020HandKey, Ji2018Local}.
    \item \textbf{Artificial Dwell Delays:} Because standard webcams lack 3D depth perception along the optical Z-axis, prior systems force users to freeze their finger over a key for 500 to 1000~ms to confirm a press, making natural touch typing impossible \cite{Khare2019QWERTY, Chen2024Gaze}.
    \item \textbf{Environmental Sensitivity:} Classical shadow and skin-color thresholding break under diffuse lighting, multiple light sources, and diverse skin tones \cite{Thomas2013Camera, Yue2014Blind}.
    \item \textbf{Rigid Mounts and Paper Shifts:} Prior systems require rigid perpendicular camera mounts; any paper displacement or camera tilt breaks key registration \cite{Wang2016AprilTag2, Kallwies2020Determining}.
    \item \textbf{Fixed Layouts:} Existing applications hard-code fixed QWERTY layouts with no action multiplexing \cite{ReviewVirtualKeyboard2020}.
\end{enumerate}

\subsection{Consequences}
\label{subsec:irca_consequences}
Consequently, monocular vision-based keyboards have failed to become viable alternatives to physical hardware. Users cannot achieve natural typing speeds, interaction is slow and fatiguing, and deployments in cost-sensitive settings (such as classrooms, public kiosks, and rural healthcare facilities) remain impractical.

\subsection{Aim}
\label{subsec:irca_aim}
To resolve this contradiction, this study aims to develop and evaluate a monocular paper-based virtual keyboard system. The system combines MediaPipe 21-joint skeletal tracking, unitless hand-length scale normalization ($L_{\text{hand}}$), lightweight PyTorch LSTM sequence classification, and dynamic AprilTag planar homography tracking ($H$) to achieve real-time ($29.09\text{ ms}$), multi-finger interaction on commodity CPUs.

\section{Research Question}
\label{sec:research_question}
To address these technical challenges, the primary research question is formulated as follows:

\begin{quote}
\textit{How can a monocular computer vision framework using strictly a single regular RGB camera and plain printed paper enable accurate, real-time, and flexible virtual keyboard interaction on flat surfaces (allowing users to design custom layouts and link different action combinations to the same physical printed sheet), executing entirely on standard commodity CPUs without requiring specialized hardware or dedicated GPUs, across regular and low camera resolutions?}
\end{quote}

This primary question is addressed through five secondary research questions (Sub-RQs):
\begin{enumerate}
    \item \textbf{Sub-RQ 1 (Fiducial Marker Selection and Layout Tracking):} Which visual fiducial marker system (such as AprilTag, ArUco, ARTag, or STag) provides the highest planar homography accuracy and lowest corner jitter under perspective tilt and low video resolutions, and how can the layout be tracked reliably when markers are partially occluded?
    \item \textbf{Sub-RQ 2 (Landmark Subsets and Kinematic Representation):} Which combination of the 21 MediaPipe hand landmarks, combined with unitless hand-length scale normalization and joint velocities, provides the most discriminative signals to identify finger touch contact without temporal smoothing delays?
    \item \textbf{Sub-RQ 3 (Deep Learning Sequence Architecture Optimization):} Which neural network architecture family (1D CNN, Attention, ResNet, BiLSTM, or LSTM) achieves the highest classification accuracy (F1-score) while maintaining near real-time inference speeds (under 3 ms) on a standard commodity CPU?
    \item \textbf{Sub-RQ 4 (Temporal Windowing and CPU Synchronization):} What temporal sliding window length, step size, and sampling rate capture impact deceleration dynamics reliably while running at a 12 FPS rate on consumer CPUs without a GPU?
    \item \textbf{Sub-RQ 5 (Layout Decoupling and Action Multiplexing):} How can physical paper layout geometry be separated from digital software semantics so that a single printed paper sheet can be dynamically bound to multiple distinct action profiles without reprinting the page?
\end{enumerate}

\section{Research Motivation}
\label{sec:research_motivation}
Several practical considerations motivate this research:
\begin{itemize}
    \item \textbf{Cost-Effective Accessibility:} By relying entirely on standard printer paper, common webcams, and commodity CPUs, the system provides affordable computing input for classrooms, developing regions, and budget workstations \cite{Khare2019QWERTY, Sreenath2021Monocular}.
    \item \textbf{Flexible and Customizable Layouts:} Users can design custom keypads for numeric data entry, gaming, or assistive typing on a single sheet of paper, and bind different software actions to the same physical page without reprinting.
    \item \textbf{Hardware Independence:} Real-time execution on standard CPUs at low-to-medium resolutions (360p, 480p, 720p) enables deployment on older or budget computers without dedicated GPUs.
    \item \textbf{Hygienic and Disposable Interfaces:} Paper layouts with AprilTag anchors can be laminated for easy sanitization or discarded after use, providing a clean input solution for medical clinics and laboratory environments \cite{ReviewVirtualKeyboard2020}.
    \item \textbf{Efficient CPU-Based Deep Learning:} Implementing a lightweight PyTorch LSTM model to detect impact dynamics on standard CPUs demonstrates an effective pattern for edge vision applications \cite{Hochreiter1997LSTM, Zhang2020MediaPipe, GilMartin2023Hand}.
\end{itemize}

\section{Research Aim}
\label{sec:research_aim}
The overarching aim of this research is to design, implement, and evaluate an accessible, customizable paper-based virtual keyboard system. The system enables users to create personalized key layouts on plain paper and multiplex the same physical sheet across multiple action profiles (such as text entry, application shortcuts, and shell scripts), running in real time on a commodity CPU using a single standard monocular RGB camera without specialized hardware or dedicated GPUs.

\section{Research Objectives Overview}
\label{sec:research_text_objectives}
To achieve this aim, the research is organized into five specific research and engineering objectives. In accordance with university dissertation guidelines, these formal general and specific objectives are detailed in Chapter~\ref{chap:objectives}.

\section{Rich Picture of the Proposed Solution}
\label{sec:rich_picture}
Figure~\ref{fig:rich_picture_diagram} illustrates the overall system architecture, divided into an offline design phase and an online runtime execution pipeline.

\begin{figure}[htbp]
\centering
\resizebox{0.65\textwidth}{!}{%
\begin{tikzpicture}[
    box/.style={draw=blue!70!black, rectangle, rounded corners=3pt,
                text width=3.0cm, minimum height=0.9cm, align=center,
                fill=blue!8!white, font=\sffamily\tiny, thick},
    process/.style={draw=teal!70!black, rectangle, rounded corners=3pt,
                text width=3.0cm, minimum height=0.9cm, align=center,
                fill=teal!8!white, font=\sffamily\tiny, thick},
    model/.style={draw=orange!80!black, rectangle, rounded corners=3pt,
                text width=3.0cm, minimum height=0.9cm, align=center,
                fill=orange!10!white, font=\sffamily\tiny, thick},
    artifact/.style={draw=purple!70!black, rectangle, rounded corners=3pt,
                text width=3.0cm, minimum height=0.9cm, align=center,
                fill=purple!8!white, font=\sffamily\tiny, thick},
    arrow/.style={-Stealth, thick, draw=gray!70}
]

%% ---- LEFT COLUMN: Offline Design Pipeline (x=0) ----
\node [box]      (designer)   at (0,  0)    {\textbf{PySide6 Layout GUI}\\Layout and Anchor Design};
\node [artifact] (xml)        at (0, -1.5)  {\textbf{Layout XML File}\\Key Bounds and Action Map};
\node [artifact] (pdf)        at (0, -3.0)  {\textbf{Printable PDF Sheet}\\Embedded AprilTag Anchors};

%% ---- RIGHT COLUMN: Online Runtime Pipeline (x=4.4) ----
\node [process]  (paper)      at (4.4, -3.0)  {\textbf{Physical Paper Setup}\\Printed Sheet on Surface};
\node [process]  (camera)     at (4.4, -4.5)  {\textbf{Monocular RGB Camera}\\Live Video Stream};
\node [box]      (apriltag)   at (4.4, -6.0)  {\textbf{AprilTag Tracker}\\Sub-Pixel Quad Detection};
\node [model]    (homography) at (4.4, -7.5)  {\textbf{Homography Matrix $H$}\\$3\times3$ Planar Warp};
\node [box]      (mediapipe)  at (4.4, -9.0)  {\textbf{MediaPipe Landmarker}\\21 Hand Joint Keypoints};
\node [process]  (norm)       at (4.4,-10.5)  {\textbf{Scale Normalisation}\\Unitless $L_{\text{hand}}$};
\node [process]  (window)     at (4.4,-12.0)  {\textbf{Temporal Windowing}\\5-Frame Sliding Window};
\node [model]    (lstm)       at (4.4,-13.5)  {\textbf{PyTorch LSTM Classifier}\\Touch Event Detection};
\node [process]  (mapping)    at (4.4,-15.0)  {\textbf{Planar Coordinate Mapping}\\$P_{\text{XML}}=H\cdot P_{\text{pixel}}$};
\node [artifact] (execution)  at (4.4,-16.5)  {\textbf{Key / Command Execution}\\OS Keystroke or Shell Script};

%% ---- LEFT COLUMN ARROWS ----
\draw [arrow] (designer) -- (xml);
\draw [arrow] (xml)      -- (pdf);

%% ---- BRIDGE: PDF to physical paper ----
\draw [arrow] (pdf.east) -- (paper.west);

%% ---- RIGHT COLUMN ARROWS ----
\draw [arrow] (paper)      -- (camera);
\draw [arrow] (camera)     -- (apriltag);
\draw [arrow] (apriltag)   -- (homography);
\draw [arrow] (homography) -- (mediapipe);
\draw [arrow] (mediapipe)  -- (norm);
\draw [arrow] (norm)       -- (window);
\draw [arrow] (window)     -- (lstm);
\draw [arrow] (lstm)       -- (mapping);
\draw [arrow] (mapping)    -- (execution);

%% ---- XML load feed ----
\draw [arrow] (xml.west) -- ++(-0.7,0) |- (execution.west);

\end{tikzpicture}%
}
\caption{Rich Picture of the Proposed Solution and Operational Pipeline Architecture.}
\label{fig:rich_picture_diagram}
\end{figure}

The pipeline executes through seven operational steps:
\begin{enumerate}
    \item \textbf{Layout Design and Export:} The user designs key boundaries and assigns commands using the PySide6 designer, exporting synchronized XML configuration and printable AprilTag PDF files.
    \item \textbf{Physical Setup:} The PDF layout is printed on standard paper, placed on a flat surface, and captured by an ordinary RGB webcam.
    \item \textbf{AprilTag Tracking and Homography:} Corner fiducials are tracked to compute the $3 \times 3$ planar homography matrix $H$, mapping image pixels to layout coordinates \cite{Wang2016AprilTag2, Kallwies2020Determining}.
    \item \textbf{Hand Tracking and Normalization:} MediaPipe extracts 21 skeletal landmarks \cite{Zhang2020MediaPipe}, which are scale-normalized relative to unitless hand length ($L_{\text{hand}}$) without temporal smoothing.
    \item \textbf{Temporal Windowing:} Normalized coordinates and joint velocities are grouped into 5-frame sliding windows with 2-frame overlap at 12 FPS \cite{Simao2016Unsupervised}.
    \item \textbf{CPU-Based LSTM Touch Detection:} A compact PyTorch LSTM network evaluates each 5-frame sequence on the CPU, classifying touch probabilities across all five fingers independently \cite{Hochreiter1997LSTM, GilMartin2023Hand}.
    \item \textbf{Coordinate Mapping and Execution:} When touch contact is confirmed, fingertip pixel coordinates are transformed via $H$ into XML coordinates ($P_{\text{XML}} = H \cdot P_{\text{pixel}}$) to dispatch the corresponding key press or system command.
\end{enumerate}

\section{Resource Requirements}
\label{sec:resource_requirements}

\subsection{Hardware Requirements}
\label{subsec:hardware_req}
\begin{itemize}
    \item \textbf{Camera Sensor:} Standard USB webcam or integrated laptop camera supporting resolutions from 360p to 1080p. High-speed cameras or depth sensors are not required.
    \item \textbf{Host Computer:} Standard commodity PC or laptop with an x86\_64 or ARM processor (such as Intel Core i3/i5 or AMD Ryzen) and at least 4 GB of RAM.
    \item \textbf{Graphics Hardware:} Dedicated GPUs are not required; all runtime inference runs entirely on the host CPU.
    \item \textbf{Physical Media:} Standard desktop laser or inkjet printer using ordinary A4 or US Letter paper.
\end{itemize}

\subsection{Software Requirements}
\label{subsec:software_req}
\begin{itemize}
    \item \textbf{Operating System:} Linux (Ubuntu 22.04/24.04 LTS), macOS, or Windows 11.
    \item \textbf{Runtime Environment:} Python 3.10+ with the \texttt{uv} package manager.
    \item \textbf{GUI Framework:} PySide6 (Qt for Python) for the visual layout designer.
    \item \textbf{Computer Vision Libraries:} OpenCV and Pupil AprilTag for image processing and tag tracking.
    \item \textbf{Pose Estimation:} Google MediaPipe Hand Landmarker for skeletal joint extraction.
    \item \textbf{Deep Learning Framework:} PyTorch for LSTM sequence model training and CPU inference.
\end{itemize}

\section{Project Scope}
\label{sec:project_scope}
Table~\ref{tab:project_scope} defines the technical boundaries of the research.

\begin{table}[H]
\centering
\caption{System Project Scope Boundaries (In-Scope vs. Out-of-Scope)}
\label{tab:project_scope}
\small
\begin{tabularx}{\textwidth}{p{0.48\textwidth} p{0.48\textwidth}}
\toprule
\textbf{In-Scope System Elements} & \textbf{Out-of-Scope System Elements} \\
\midrule
Single standard monocular RGB camera (USB or laptop webcam) operating at resolutions of 360p, 480p, 720p, or 1080p at 12 FPS. & Depth cameras, Time-of-Flight sensors, infrared illuminators, or stereo camera rigs. \\
\addlinespace
Single active hand tracking with 21 landmarks, evaluating independent touch probabilities across all five fingers. & Bimanual (two-handed) typing, deferred to future work to avoid tag occlusion on compact A4 paper. \\
\addlinespace
Standard plain paper with four border AprilTag anchors printed via desktop printer (A4 or Letter). & Active capacitive overlays, conductive ink, wired panels, or electronic paper components. \\
\addlinespace
Real-time CPU inference on commodity processors without requiring dedicated GPUs. & High-end GPU workstations or cloud-based inference servers. \\
\addlinespace
PySide6 desktop application for designing custom layouts, exporting XML maps, and generating printable PDFs. & Mobile operating system apps (Android/iOS) or browser extensions. \\
\addlinespace
AprilTag fiducial tracking with $3 \times 3$ planar homography under perspective tilt. & Non-planar or curved 3D surfaces (such as cylinders or spheres). \\
\addlinespace
Unitless hand-length scale normalization ($L_{\text{hand}}$) with direct feature propagation for resolution invariance. & Biomechanical muscle force modeling or finger joint pressure estimation. \\
\addlinespace
Compact PyTorch LSTM sequence classifier using 5-frame sliding windows on CPU. & Large language model auto-correction or word prediction engines. \\
\addlinespace
Real-time keystroke dispatch and XML key boundary lookup for flat surfaces. & Multi-user simultaneous input on a single paper sheet. \\
\bottomrule
\end{tabularx}
\end{table}

\subsection{In-Scope Technical Boundaries}
\label{subsec:in_scope_boundaries}
The system operates using a single standard monocular RGB webcam and plain printed paper \cite{Posner2012Single, Enkhbat2020HandKey, Sreenath2021Monocular}. It functions reliably on common low-to-medium resolutions (360p, 480p, 720p, 1080p) without requiring studio optics or depth sensors. Layouts are printed on standard A4 or Letter paper with no embedded electronics.

Interaction is optimized for single-handed typing, tracking 21 skeletal joints of one active hand. Touch probabilities are evaluated independently across all five fingers (Thumb, Index, Middle, Ring, Pinky). Single-handed tracking achieves a low latency of 29.09 ms on standard CPUs and prevents hand occlusion of the border AprilTags on compact A4 sheets. The PySide6 Layout Designer decouples physical geometry from digital semantics, allowing one printed sheet to be bound to multiple action profiles (text typing, developer shortcuts, audio controls, shell commands) without reprinting \cite{Doan2022Efficient, Li2021CNN, Lee2022Virtual}. AprilTag homography dynamically compensates for camera tilt and paper movement \cite{Wang2016AprilTag2, Kalaitzakis2021Fiducial}. Landmark coordinates are scale-normalized relative to hand length ($L_{\text{hand}}$) at 12 FPS and classified using the lightweight PyTorch LSTM model on standard CPUs \cite{Hochreiter1997LSTM, Zhang2020MediaPipe, GilMartin2023Hand}.

\subsection{Out-of-Scope Boundaries and Future Extensions}
\label{subsec:out_scope_boundaries}
To ensure feasibility and low-cost deployment on commodity hardware, several elements remain outside the scope of this project. Specialized hardware sensors (such as ToF cameras, stereo pairs, infrared projectors, and data gloves) are excluded \cite{Lee2019Virtual, Kudale2016RealTime}. Non-planar 3D surfaces (such as cylinders or spheres) are omitted, focusing strictly on flat desktop planes. Two-handed bimanual typing and multi-user interaction are excluded from the current prototype to prevent AprilTag occlusion on standard A4 paper and minimize CPU load. Finally, physical touch force measurement and language model predictive auto-correction are reserved for future work.

\section{Thesis Organization}
\label{sec:thesis_organization}
The remainder of this dissertation is organized into five subsequent chapters:
\begin{itemize}
    \item \textbf{Chapter 2 (Objectives):} Defines the general research aim and details the five specific research and engineering objectives driving this work.
    \item \textbf{Chapter 3 (Literature Review):} Evaluates prior work in spatial HCI, optical sensing modalities, fiducial planar homography, and deep learning touch detection, synthesizing existing research gaps.
    \item \textbf{Chapter 4 (Methodology):} Details the design science research methodology, the two-application software architecture, the 12 FPS kinematic preprocessing pipeline, scale normalization, and PyTorch deep learning models.
    \item \textbf{Chapter 5 (Results):} Presents empirical benchmarking results for deep learning sequence classifiers, landmark ablations, planar homography precision, CPU execution latency, and typing usability.
    \item \textbf{Chapter 6 (Discussion and Conclusions):} Analyzes experimental findings, evaluates system limitations, highlights practical application domains, outlines future research extensions, and synthesizes final conclusions.
\end{itemize}

\section{Chapter Summary}
\label{sec:chapter_summary}
This introductory chapter presented the research problem background, formal problem statement, primary research question and sub-questions, practical motivations, high-level operational pipeline, resource requirements, and project scope boundaries. By relying strictly on a single standard monocular RGB camera, plain printed paper with AprilTag anchors, and lightweight deep learning running on a commodity CPU, the framework addresses longstanding barriers in accessible human-computer interaction. The formal research objectives guiding the execution of this dissertation are detailed in Chapter~\ref{chap:objectives}.
